\chapter{Introduction}
\label{sec:introduction}

%\HC{Perhaps cite~\cite{10.1002_cpe.8019} as another example of a real-time workflow? Not sure}

Executions on compute platforms are subject to hardware
and software errors. An active area of research is application scheduling
in the presence of failures in non-real-time~\cite{10.5555/2811302}
or real-time~\cite{10.1145/3589950} settings, typically using
rollback/recovery and/or computation replication techniques.  In this work, 
we focus on deadline-constrained scientific computing applications, which
are increasingly commonplace across the sciences (see a recent
U.S. Department of Energy study~\cite{IRI-workflows}).

In many domains, approximate computing techniques~\cite{10.1145/3716845}
make it possible to implement application tasks with multiple \emph{execution
modes}. Each mode corresponds to a trade-off between execution
time and output error, where an erroneous output retains some utility.
Trading off output quality for a
shorter execution time provides an opportunity to meet real-time deadlines
that would otherwise be missed, motivating the development of strategies
for scheduling approximate computing applications in the presence of
failures~\cite{10.1145/2442087.2442094, 10.1007/s11227-021-03764-x,
259428, 10.1177/0037549709340729, 10.1109/TC.2003.1183950, 7570524}.


In this work we consider an application structured as a chain of tasks,
some of which have multiple modes of execution. The application executes on
a platform comprised of one or more processing elements (PEs) and is
subject to fail-stop failures.  Given an end-to-end deadline, the objective
is to make scheduling decisions that minimize the expected error of the
best final output produced by that deadline, assigning a user-defined penalty
when no final output is produced.  Prior work~\cite{10678736}
considers static replication and resource provisioning in this setting. 
Instead, we consider dynamic scheduling decisions 
that can employ replication, rollback/recovery,
execution mode changes after failures, reuse of persistent intermediate
outputs, and coordinated scheduling across PEs.  Specifically, we make the
following contributions:
\begin{compactitem}
    \item We formulate the problem and show its NP-hardness;
    \item For one-PE platforms we discretize time to develop a dynamic programming approach, together with a practical method for selecting the discretization resolution;
    \item For multiple-PE platforms we propose heuristics, some of which build on our results for the one-PE case; and
    \item We evaluate our scheduling strategies in simulation, both for a production workload from the field of fusion energy science and for a large set of synthetic workloads. 
\end{compactitem}

The rest of this thesis is organized as follows. Chapter~\ref{sec:use-case}
details a production use case that motivates our scheduling problem, which
is formalized in Chapter~\ref{sec:problem}.  Chapter~\ref{sec:methodology}
describes our experimental methodology. Chapters~\ref{sec:onehost} and
\ref{sec:multihost} present our approach and results for one-PE and
multiple-PE platforms, respectively. 
Chapter~\ref{sec:related} reviews related work.  Finally,
Chapter~\ref{sec:conclusion} concludes with a brief summary of results and
perspectives on future work.
